% =============================================================================
% CHƯƠNG 5 — KẾT LUẬN VÀ HÀM Ý
% Văn phong khoa học, thuật ngữ và cấu trúc đã được chuẩn hóa.
% =============================================================================

\chapter{KẾT LUẬN VÀ HÀM Ý}
\label{chap:ket-luan}

\section{Kết luận chính}
\label{sec:c5-main-conclusion}

Nghiên cứu so sánh GBM đa biến, Heston đa biến và
\textit{Schrödinger Bridge Time Series} trong cùng một quy trình
\textit{Deep Hedging}. Hai hợp đồng gồm basket Asian call và Asian worst-of
put trên rổ AAPL--JPM--XOM. Mỗi hợp đồng được đánh giá tại ba mức giá thực
hiện. Kiến trúc mạng, số quỹ đạo, chi phí giao dịch và quy trình MSE--CVaR được
giữ cố định giữa ba mô hình sinh dữ liệu.

Trong giai đoạn đại diện (2023--2025), SBTS có độ lệch chuẩn sai số hedging thấp nhất tại
sáu tổ hợp, giảm 13--42\% so với GBM và 8--38\% so với Heston. MCS giữ SBTS là
phần tử duy nhất tại cả sáu tổ hợp. Trong giai đoạn phục hồi (2021--2022), SBTS được lựa
chọn cho basket call nhưng không được lựa chọn đồng nhất cho worst-of put.
Trong giai đoạn căng thẳng (2019--2020), mô hình tham số có $\sigma(R)$ thấp hơn SBTS tại năm
trong sáu tổ hợp của mỗi cặp so sánh; MCS giữ cả ba mô hình tại năm tổ hợp.

SBTS có $V_0^{*}$ thấp nhất tại sáu tổ hợp, giảm 13--45\% so với GBM và
17--44\% so với Heston. Trong giai đoạn MSE, $V_0^{*}$ là mức vốn ban đầu làm
sai số hedging kỳ vọng bằng không dưới phân phối huấn luyện. Đại lượng này
còn được gọi là giá bàng quan bậc hai (\textit{quadratic indifference price}).
Thứ hạng $V_0^{*}$ được diễn giải cùng thứ hạng sai số hedging tại từng giai
đoạn thị trường.

\section{Trả lời câu hỏi nghiên cứu}
\label{sec:c5-rq}

\paragraph{RQ1 -- Mô hình sinh dữ liệu dựa trên dữ liệu có làm giảm sai số hedging?}
SBTS làm giảm $\sigma(R)$ tại sáu tổ hợp của giai đoạn đại diện và ba tổ hợp basket call
của giai đoạn phục hồi. Kết quả không được duy trì tại worst-of put của giai
đoạn phục hồi và tại giai đoạn căng thẳng. Hiệu
quả phụ thuộc vào giai đoạn thị trường và loại hợp đồng.

\paragraph{RQ2 -- $V_0^{*}$ có khác biệt hệ thống giữa các mô hình sinh dữ liệu?}
SBTS có $V_0^{*}$ thấp nhất tại sáu tổ hợp. Chênh lệch phản ánh phân phối khoản
thanh toán, P\&L và chi phí giao dịch dưới từng mô hình sinh dữ liệu. Kết quả
so sánh mức vốn ban đầu theo tiêu chuẩn bình phương sai số, không phải so sánh giá
thị trường hoặc giá phi arbitrage.

\paragraph{RQ3 -- Thứ hạng có ổn định qua các giai đoạn thị trường?}
Thứ hạng không ổn định. Giai đoạn đại diện lựa chọn SBTS. Giai đoạn phục hồi
cho kết quả khác nhau giữa hai hợp đồng. Trong giai đoạn căng thẳng, thứ hạng
đảo chiều ở phần lớn kiểm định từng cặp và MCS không chọn một mô hình duy nhất
tại năm tổ hợp.

\section{Đánh giá các giả thuyết}
\label{sec:c5-hypotheses}

\begin{table}[H]
\centering
\small
\caption{Tổng hợp đánh giá giả thuyết nghiên cứu}
\label{tab:c5-hypotheses}
\begin{tabularx}{\textwidth}{@{}
>{\raggedright\arraybackslash}p{0.09\textwidth}
>{\raggedright\arraybackslash}p{0.25\textwidth}
X@{}}
\toprule
\textbf{Giả thuyết} & \textbf{Kết luận} & \textbf{Bằng chứng chính}\\
\midrule
H1 & Được ủng hộ một phần & SBTS có $\sigma(R)$ thấp nhất tại sáu tổ hợp của giai đoạn
đại diện. Sáu so sánh với GBM và năm trong sáu so sánh với Heston có
$p_{\mathrm{BH}}<0{,}05$.\\
H2 & Được ủng hộ về mô tả & $V_0^{*}$ của SBTS thấp nhất tại sáu tổ hợp, giảm
13--45\% so với GBM và 17--44\% so với Heston. Kết quả $V_0^{*}$ không có kiểm
định riêng.\\
H3 & Được ủng hộ & MCS ở mức ý nghĩa 10\% chỉ giữ SBTS tại sáu tổ hợp của giai đoạn đại diện.\\
H4 & Được ủng hộ về thứ hạng & Năm trong sáu tổ hợp của giai đoạn căng thẳng của mỗi cặp
SBTS--mô hình tham số có $p_{\mathrm{BH}}<0{,}05$ theo hướng bất lợi cho SBTS.
Cơ chế ngoại suy của kernel chưa được kiểm định riêng.\\
\bottomrule
\end{tabularx}
\end{table}

Đánh giá giả thuyết phân biệt kết quả thực nghiệm và cơ chế giải thích. Bộ kết
quả xác định sự thay đổi thứ hạng giữa các giai đoạn thị trường nhưng không có
kiểm định riêng về mức co của hỗ trợ kernel.

\section{Đóng góp của nghiên cứu}
\label{sec:c5-contributions}

Thứ nhất, nghiên cứu so sánh ba mô hình sinh dữ liệu đa biến trong cùng quy
trình \textit{Deep Hedging} cho quyền chọn Asian rainbow dưới chi phí giao
dịch. Kiến trúc mạng, hợp đồng, mức giá thực hiện, số quỹ đạo và quy trình tối
ưu được giữ cố định.

Thứ hai, nghiên cứu phân tích $V_0^{*}$ như mức vốn ban đầu theo tiêu chuẩn
bình phương sai số. Kết quả cung cấp đồng thời thứ hạng về sai số hedging và
thứ hạng về mức vốn ban đầu.

Thứ ba, nghiên cứu đánh giá thứ hạng tại ba giai đoạn thị trường bằng kiểm định
từng cặp và MCS. Kết quả tách biệt thứ hạng tại giai đoạn đại diện, phục hồi và căng thẳng.

\section{Hàm ý thực tiễn}
\label{sec:c5-implications}

Việc lựa chọn mô hình sinh dữ liệu cần sử dụng kết quả hedging rủi ro trên
dữ liệu đánh giá, bên cạnh các chỉ tiêu phân phối. Kết quả phải được tách theo
hợp đồng và giai đoạn thị trường. Trong giai đoạn đại diện, SBTS có $\sigma(R)$ và
$V_0^{*}$ thấp hơn hai mô hình tham số. Trong giai đoạn căng thẳng, thứ hạng $\sigma(R)$ đảo
chiều và MCS không chọn một mô hình duy nhất tại năm trong sáu tổ hợp.

Checkpoint sau MSE và checkpoint sau CVaR cần được đánh giá riêng. Tại
$\kappa=1{,}00$ trong giai đoạn đại diện, giai đoạn CVaR làm $\sigma(R)$ của basket call
SBTS giảm 19,3\% nhưng làm $\sigma(R)$ của worst-of put SBTS tăng 13,4\%.

\section{Hạn chế của nghiên cứu}
\label{sec:c5-limitations}

\paragraph{Giới hạn dữ liệu và số chiều.}
Nghiên cứu chỉ dùng ba cổ phiếu Hoa Kỳ và một cửa sổ hiệu chỉnh 2005--2019.
Kết quả không được suy rộng sang rổ có số chiều lớn hơn, ngoại hối hoặc tài
sản số.

\paragraph{Giới hạn hợp đồng.}
Nghiên cứu chỉ xem xét basket Asian call và Asian worst-of put. Quyền chọn kiểu
Mỹ, Bermudan, best-of, digital rainbow, barrier và autocallable không thuộc
phạm vi thực nghiệm.

\paragraph{Đặc tả Heston giản lược.}
Heston dùng tham số động lực chung và không hiệu chỉnh từ bề mặt biến động hàm
ý theo từng tài sản. Kết quả chỉ áp dụng cho đặc tả Heston tại
Chương~\ref{chap:phuong-phap} \parencite{Dimitroff2011}.

\paragraph{Một kiến trúc và một lịch huấn luyện.}
Kiến trúc $10\to64\to64\to3$ và lịch 500+200 epoch được giữ cố định. LSTM,
Transformer, mạng có vùng không giao dịch và lịch huấn luyện thích nghi không
được so sánh.

\paragraph{Chồng lấn lịch và hiệu chỉnh tĩnh.}
Năm 2019 vừa thuộc cửa sổ hiệu chỉnh mô hình sinh dữ liệu vừa thuộc giai đoạn căng thẳng
2019--2020. Mạng không học trực tiếp dữ liệu lịch sử, nhưng giai đoạn này không
phải ngoài mẫu hoàn toàn đối với bước hiệu chỉnh mô hình sinh dữ liệu. Ba mô
hình được hiệu chỉnh một lần và không được cập nhật theo cửa sổ cuốn.

\paragraph{Seed thay thế và tính ghép cặp.}
Checkpoint SBTS--worst-of put--$\kappa=0{,}95$--seed 3 phân kỳ và được thay
bằng seed 10. Tổ hợp này có mười lần chạy hội tụ nhưng một quan sát không ghép
cùng seed với GBM và Heston. Hạn chế này áp dụng cho các kiểm định liên quan
đến tổ hợp trên.

\paragraph{Giới hạn của $V_0^{*}$.}
$V_0^{*}$ là nghiệm của bài toán hedging theo bình phương sai số dưới phân
phối của mô hình sinh dữ liệu. Đại lượng này không phải giá trung hòa rủi ro,
không phải giá thị trường và không bao gồm các yêu cầu vốn pháp định.

\paragraph{Giới hạn nhận dạng cơ chế stress.}
Kết quả xác định sự đảo chiều trong giai đoạn căng thẳng nhưng không đo trực tiếp số lân cận
kernel hoặc khoảng cách từ trạng thái đánh giá đến tập tham chiếu. Cơ chế suy
giảm ngoài miền dữ liệu chưa được kiểm định độc lập.

\section{Hướng nghiên cứu tiếp theo}
\label{sec:c5-future}

Hướng thứ nhất là tái lập toàn bộ lưới với tập seed chung được xác định trước,
đồng thời bổ sung bootstrap và kiểm định hoán vị cho mẫu nhỏ.

Hướng thứ hai là đo số lân cận kernel có trọng số khác không, khoảng cách từ
trạng thái đến tập tham chiếu và sai số hedging theo mức dịch chuyển phân
phối.

Hướng thứ ba là cập nhật SBTS theo cửa sổ cuốn, bổ sung dữ liệu căng thẳng và
tính lại MCS trên cửa sổ đánh giá gần nhất.

Hướng thứ tư là sử dụng quá trình tham chiếu có bước nhảy hoặc bổ sung các kịch
bản căng thẳng có kiểm soát. Thực nghiệm cần giữ cùng thiết kế hợp đồng--mức
giá thực hiện--seed giữa các mô hình.

Hướng thứ năm là mở rộng sang autocallable worst-of, barrier rainbow, rổ nhiều
tài sản, LSTM, Transformer và mạng có vùng không giao dịch.

\section{Kết luận chương}

Trong giai đoạn đại diện, SBTS có $\sigma(R)$ và $V_0^{*}$ thấp hơn hai mô hình tham số;
MCS chỉ giữ SBTS tại sáu tổ hợp. Trong giai đoạn phục hồi, kết quả phụ thuộc vào loại
hợp đồng. Trong giai đoạn căng thẳng, thứ hạng sai số hedging đảo chiều và MCS giữ nhiều mô
hình tại năm trong sáu tổ hợp. Do đó, kết luận được xác định theo mô hình sinh
dữ liệu, hợp đồng, hàm mục tiêu và giai đoạn thị trường.
